\section{Conclusiones}

A lo largo de este trabajo práctico, logramos analizar y clasificar la complejidad del problema de \textit{Hitting-Set}. En primer lugar, construimos un algoritmo validador que nos permitió confirmar que el problema pertenece a la clase NP. En palabras simples, esto significa que si alguien nos propone una lista de jugadores como solución, podemos verificar rápidamente (en tiempo polinomial) si esa selección cumple con las exigencias de todos los medios de prensa.

Posteriormente, demostramos que el problema es NP-Completo utilizando una técnica conocida como reducción, partiendo del problema de \textit{Dominating Set}. La utilidad de aplicar reducciones en la computación es enorme: nos permite usar lo que ya sabemos para clasificar problemas nuevos. Como ya está demostrado que \textit{Dominating Set} es un problema computacionalmente "difícil", al lograr transformarlo en nuestro problema de \textit{Hitting-Set}, probamos que el nuestro hereda esa misma dificultad. Básicamente, esta técnica nos ahorra el esfuerzo de tener que analizar la dificultad del problema desde cero, dándonos una garantía matemática sobre a qué nos enfrentamos.

Llevando todo este análisis teórico al caso real (la necesidad del cuerpo técnico de seleccionar jugadores para contentar a la prensa), haber comprobado que este problema es NP-Completo tiene una consecuencia directa: es altamente improbable que exista una "fórmula mágica" o un algoritmo rápido que nos dé la lista óptima de jugadores en poco tiempo para cualquier cantidad de datos.

Saber esto justifica formalmente por qué debemos cambiar nuestra estrategia de programación en las siguientes etapas del trabajo. Para listas de convocados de tamaño moderado, no nos queda otra opción que usar métodos de búsqueda exhaustiva inteligentes, como el \textit{Backtracking}, que analizan muchas combinaciones pero descartan rápidamente las que no sirven. Por otro lado, si en el futuro la cantidad de jugadores y de periodistas creciera muchísimo, la teoría nos avisa desde ahora que lo mejor será dejar de buscar el equipo perfecto y empezar a usar heurísticas o algoritmos de aproximación, que nos entreguen un equipo "lo suficientemente bueno" en un tiempo razonable.